\documentclass[24pt,a4paper]{article}% 文档格式
\usepackage{ctex,hyperref}% 输出汉字
\usepackage{times}% 英文使用Times New Roman
%%%%%%%%%%%%%%%%%%%%%%%%%%%%%%%%%%%%%%%%%%%%%%%%%%%%%%%%
\title{\fontsize{18pt}{27pt}\selectfont% 小四字号，1.5倍行距
	{\heiti% 黑体 
		一种\LaTeX 模板}}% 题目
%%%%%%%%%%%%%%%%%%%%%%%%%%%%%%%%%%%%%%%%%%%%%%%%%%%%%%%%
\date{}% 日期（这里避免生成日期）
%%%%%%%%%%%%%%%%%%%%%%%%%%%%%%%%%%%%%%%%%%%%%%%%%%%%%%%%
\usepackage{amsmath,amsfonts,amssymb}% 为公式输入创造条件的宏包
%%%%%%%%%%%%%%%%%%%%%%%%%%%%%%%%%%%%%%%%%%%%%%%%%%%%%%%%
\usepackage{graphicx}% 图片插入宏包
\usepackage{subfigure}% 并排子图
\usepackage{float}% 浮动环境，用于调整图片位置
\usepackage[export]{adjustbox}% 防止过宽的图片
\usepackage{caption}
%%%%%%%%%%%%%%%%%%%%%%%%%%%%%%%%%%%%%%%%%%%%%%%%%%%%%%%%
\usepackage{url}% 超链接
\usepackage{bm}% 加粗部分公式
\usepackage{listings}
\usepackage{xcolor}
\lstset{
    language=Java,
    basicstyle=\ttfamily,
    keywordstyle=\color{blue},
    commentstyle=\color{gray},
    stringstyle=\color{orange}
}
\usepackage{multirow}
\usepackage{booktabs}
\usepackage{epstopdf}
\usepackage{epsfig}
\usepackage{longtable}% 长表格
\usepackage{supertabular}% 跨页表格
\usepackage{algorithm}
\usepackage{algorithmic}
\usepackage{changepage}% 换页
\usepackage{listings}% 插入代码段
%%%%%%%%%%%%%%%%%%%%%%%%%%%%%%%%%%%%%%%%%%%%%%%%%%%%%%%%
\usepackage[left=2.50cm,right=2.50cm,top=2.80cm,bottom=2.50cm]{geometry}% 页边距设置
\renewcommand{\baselinestretch}{1.5}% 定义行间距（1.5）
\renewcommand{\contentsname}{\normalfont \kaishu \Huge 目录}% 定义目录两字的格式

\usepackage{subfigure}% 有关设置目录引导点的宏包
\usepackage[subfigure]{tocloft}
\renewcommand{\cftsecleader}{\cftdotfill{\cftdotsep}} % 给 sections加点
\newcommand\mydot[1]{\scalebox{#1}{.}}
\renewcommand\cftdot{\mydot{0.8}} % change the size of dots
\renewcommand\cftdotsep{3} % change the space between dots

\hypersetup{
colorlinks=true,
linkcolor=black
}% 设置链接的颜色，防止目录出现红框
%%%%%%%%%%%%%%%%%%%%%%%%%%%%%%%%%%%%%%%%%%%%%%%%%%%%%%%%

\begin{document}% 以下为正文内容
\begin{titlepage}
    \centering
    \includegraphics[width=0.3\textwidth]{image/zju_logo.png}\par\vspace{5cm}
    {\huge\songti 设计模式报告\par}
    \vspace{1cm}
    {\Large\itshape 第二小组\par}
    \vspace{7cm}

    \vfill
    {\large \today\par}
\end{titlepage}
\newpage

\begin{center}
    \kaishu
    \tableofcontents
    \setcounter{page}{0}
    \thispagestyle{empty} % 设置目录页的页脚为空
\end{center}
\newpage

\section*{\songti 第一部分：外观模式}
\addcontentsline{toc}{section}{1.外观模式}
\subsection*{\songti 1.1概念}
\addcontentsline{toc}{subsection}{1.1概念}
外观模式是一种常用的软件设计模式，属于结构型模式的一种。它为一个
子系统中的一组接口提供一个简洁的，统一的界面，用来掩盖系统的复杂
性，使得子系统中的一个接口更加容易使用。这可以通过定义一个包含子
系统所有可能功能的单一类来实现的。
\subsection*{\songti 1.2意图}
\addcontentsline{toc}{subsection}{1.2意图}
外观模式的主要意图为子系统中一组接口提供一个统一的界面，从而降低
客户端与复杂子系统之间的交互难度，通过创建统一的界面，来控制客户
端和复杂系统的交互。其核心目的在于隐藏子系统的实现细节，让子系统
使用起来更简便，从而减轻客户端的使用负担。

在实际的应用中，一个复杂系统经常拥有大量的组件与底层的模块。如果
客户端直接与这些复杂组件和模块交互，将会导致代码耦合度高，维护困
难。在这个时候，使用外观模式将这些复杂的操作封装起来，变得尤为重
要，因为通过定义一个高层、简单易用的接口库，为客户端提供一个统一
的界面，使客户端能够以一种更加理解和更易使用的方式与子系统进行交
流和操作。这个模式把复杂的业务逻辑和操作封装在一个简洁易用的“外观”
接口之后，并通过系统内部的协调和调度来完成实际的业务工作。

这样不仅降低了子系统的复杂度，也降低了客户端使用子系统的学习成
本，让客户端只需要学习接口的定义和功能即可；同时外观模式也在一定
程度上提升了代码的可阅读性和可维护性。在子系统内部需要进行变化或
升级时，客户端的代码不需要进行修改，由外观接口内部的子系统来完成
这些更改升级即可。

总体而言，外观模式的意图是将一个复杂系统背后的许多复杂交互行为隐
藏起来，面向客户端来提供一个简单易用、一致性的界面，从而减轻客户
端的使用复杂度，使子系统的使用更加方便。
\subsection*{\songti 1.3结构和组成}
\addcontentsline{toc}{subsection}{1.3结构和组成}
外观模式主要由以下几部分组成：
\subsubsection*{\songti 1.3.1\ 外观类}
外观类是外观模式的核心，客户端通过调用该类来访问子系统，该类
知道每个子系统的功能和责任，将客户端的请求代理给适当的子系统
对象。
\subsubsection*{\songti 1.3.2\ 子系统类}
在外观模式中，可以有一个或多个子系统。每个子系统并
不是单一的类，而是由多个类组成。在子系统里，每个类都有各自的
功能，这些功能之间可以相互交互，共同协作完成任务。但对外部客
户端来说，这些交互被隐藏起来不可见，客户端只可以通过外观类来
访问子系统的各个功能。
\subsubsection*{\songti 1.3.3\ 客户端}
在外观模式中，客户端通过调用外观类中的功能来完成对于
整个子系统的访问，而无需关心子系统的内部工作细节。
\subsection*{\songti 1.4优点和缺点}
\addcontentsline{toc}{subsection}{1.4优点和缺点}

\subsubsection*{\songti 1.4.1\ 优点}
\begin{itemize}
    \item \textbf{简化客户端的复杂度}：外观模式隐藏了系统的复杂性，使客户端代码
更加简洁和易于使用。
    \item \textbf{解耦系统各个部分}：外观模式可以将系统各个部分解耦，使它们能够
独立变化。这样，系统的某个部分的变化不会影响到其他部分，从而
便于各个部分的更改。
    \item \textbf{提高系统的可维护性}：外观模式将系统的复杂性隐藏在了一个单独的
外观类中，使代码更加清晰和易于维护。
    \item \textbf{提高系统的安全性}：外观模式可以控制客户端能够访问系统的部分，
从而提高系统的安全性。
\end{itemize}

\subsubsection*{\songti 1.4.2\ 缺点}
\begin{itemize}
    \item \textbf{缺乏细粒度控制}：不外观模式提供了一组简单的接口，用
于访问系统的各个部分。但是如果需要更细粒度的控制，可能需要修
改外观类或者直接访问系统的各个部分。
    \item \textbf{增加复杂性}：外观模式需要引入一个额外的外观类，如果客
户端需要与子系统进行大量交互，可能会增加系统的复杂性。
    \item \textbf{影响性能}如果外观类的实现不够高效，可能会影响系
统的性能。
    \item \textbf{提高系统的安全性}：外观模式可以控制客户端能够访问系统的部分，
从而提高系统的安全性。
\end{itemize}

\subsection*{\songti 1.5应用场景}
\addcontentsline{toc}{subsection}{1.5应用场景}
\begin{enumerate}
    \item 设计一个简单的方法以便于客户端调用复杂的系统。例如，为复杂的
计算机系统设计一个启动按钮。在客户按下启动按钮后，计算机开始
进行内部检查、完成系统的加载、进行进程的调度等一系列操作，而
对于客户来说，仅需要点击一个按钮就可以了，而不用关心后面的复
杂过程。
    \item 当需要定义系统的多级访问权限时，可以通过在系统一级接口的基础
上封装一层访问接口，控制用户对于某些高级接口的访问。例如在网
络服务器的设计中，在不同级别的用户接口上再封装一层外观接口，
从而控制用户对于网络资源的访问。
    \item 将子系统和客户端解耦，使得他们之间的依赖关系降低，同时也使子
系统更易于独立开发和维护，适合在大型系统设计时使用。
\end{enumerate}

\subsection*{\songti 1.6工业界的实际应用}
\addcontentsline{toc}{subsection}{1.6工业界的实际应用}
\begin{enumerate}
    \item \textbf{Java的Spring框架}：Spring框架提供了一个统一的架构，把复杂的事
务管理、安全性、数据库访问等功能封装成为了一个统一的接口，使
得用户可以非常方便地使用这些复杂功能。
    \item \textbf{数据库连接}：许多应用程序使用的数据库连接库，例如 Java 的
JDBC，为复杂的 SQL 查询和数据库操作提供了简单的接口。通过这
些接口，程序员不再需要直接编写 SQL 语句，只需要调用一些易于理
解的方法和函数就能进行数据库的操作。
    \item \textbf{电商平台 API}：例如 Amazon、eBay 和 阿里巴巴 这样的电商平台，
它们提供了一系列的API接口来帮助开发者进行商品上传、查询、售卖
和物流管理等一系列操作。商家不需要关注具体的实现细节，只需要
调用相应的接口就可以方便地进行交易。
    \item \textbf{GUI库}：在许多 GUI（图形用户界面）库中，例如Swing、Qt等，都使
用了外观模式思想，他们为复杂的底层图形系统提供了简单易用的API
接口，牵扯到窗口管理、事件处理等一系列复杂的操作，却只需要调
用简单的接口函数。这大幅度降低了GUI编程难度，提高开发效率。
\end{enumerate}

\subsection*{\songti 1.7本项目中的应用}
\addcontentsline{toc}{subsection}{1.7本项目中的应用}
本项目中，外观类主要运用在工具类和底层数据的使用上。对于一些需要频繁使用的工具类和算法，例如 ThreadLocal、Mail 和 MD5，本项目的后端选择将其集成为工具类，供其他类使用。这样可以大大简化系统代码，并很好地解耦了高级类和低级方法。

一个例子如下
\begin{lstlisting}[language=Java]
@SuppressWarnings("all")
public class ThreadLocalUtil {
    // 提供ThreadLocal对象,
    private static final ThreadLocal THREAD_LOCAL = new ThreadLocal();
    // 根据键获取值
    public static <T> T get() {
        return (T) THREAD_LOCAL.get();
    }
    // 存储键值对
    public static void set(Object value) {
        THREAD_LOCAL.set(value);
    }
    // 清除ThreadLocal 防止内存泄漏
    public static void remove() {
        THREAD_LOCAL.remove();
    }
}
\end{lstlisting}
我们将提供线程局部变量的 ThreadLocal 包装为一个外观类，从而简化了其他类对其的使用。

此外，我们还将针对数据库的操作使用外观模式包装，将数据库操作与业务逻辑分离，使得业务逻辑代码更加清晰和可维护。

一个例子如下：
\begin{lstlisting}[language=Java]
@Mapper
public interface AccountMapper {
    @Update("update account set balance=#{balance} where acno=#{acno}")
    void updateAccount(int acno, double balance);
    
    @Select("select * from account where pano=#{pano}")
    Account findAccountByPano(int pano);
    
    @Select("select Balance from account where pano=#{pano}")
    double findBalanceByPano(int pano);
}
\end{lstlisting}
我们将一些必要的数据库操作包装成一个外观类，在外观类内部通过 @Select 注解指定对应的 SQL 查询语句，实现了对数据库表的操作，并将这些操作作为接口提供给其他类使用。

\newpage
\end{document}% 结束文档编辑，后面写啥都编译不出来

%%%%%%%%%%%%%%%%%%%%%草稿%%%%%%%%%%%%%%%%%%%%%
% 插入图片 %
\begin{figure}[H]
    \centering
    \includegraphics[width=1\textwidth]{images/*.png}
    \caption*{图*.键入标题...}
\end{figure}
%%%%%%%%%%%%%%%%%%%%%%%%%%%%%%%%%%%%%%%%%%%%%